% =========================================================
% 3.12 Decoder
% FINAL VERSION
% Fixed 3-column layout:
% Truth Table | Boolean Equations | Logic Circuit
% =========================================================

\section{ডিকোডার (\textenglish{Decoder})}

% =========================================================
% 3.12.1 Concept
% =========================================================

\subsection{ডিকোডারের ধারণা
(\textenglish{Concept of Decoder})}

কম্পিউটার ও ডিজিটাল সিস্টেমে তথ্য সাধারণত বাইনারি কোডে প্রক্রিয়াকরণ করা হয়।
অনেক ক্ষেত্রে এই বাইনারি কোডের ভিত্তিতে একটি নির্দিষ্ট
\textenglish{output line}, ডিভাইস বা মেমোরি অবস্থান নির্বাচন করতে হয়।

যে প্রক্রিয়ায় বাইনারি কোড অনুযায়ী নির্দিষ্ট
\textenglish{output} নির্বাচন করা হয়, তাকে
\textbf{ডিকোডিং (\textenglish{Decoding})} বলে।
এ কাজের জন্য ব্যবহৃত ডিজিটাল সার্কিটকে
\textbf{ডিকোডার (\textenglish{Decoder})} বলে।

\begin{center}
\begin{chThreeDiagram}
\begin{tikzpicture}[
    >=Stealth,
    line width=0.8pt,
    node distance=1.0cm
]

\node[
    draw,
    rounded corners=2pt,
    fill=ICTLightBlue,
    minimum width=2.8cm,
    minimum height=0.9cm
] (in) {\textenglish{Binary Input}};

\node[
    draw,
    rounded corners=2pt,
    fill=ICTLightGreen,
    minimum width=2.5cm,
    minimum height=0.9cm,
    right=of in
] (dec) {\textbf{\textenglish{Decoder}}};

\node[
    draw,
    rounded corners=2pt,
    fill=ICTLightOrange,
    minimum width=3.0cm,
    minimum height=0.9cm,
    right=of dec
] (out) {\textenglish{Selected Output}};

\draw[->] (in) -- (dec);
\draw[->] (dec) -- (out);

\end{tikzpicture}
\end{chThreeDiagram}
\end{center}

% =========================================================
% 3.12.2 Definition
% =========================================================

\subsection{ডিকোডারের সংজ্ঞা
(\textenglish{Definition of Decoder})}

\par\begingroup

যে সমবায় লজিক সার্কিট
(\textenglish{Combinational Logic Circuit})
\(n\)-বিট বাইনারি \textenglish{input} গ্রহণ করে এবং
\textenglish{input}-এর মান অনুযায়ী সর্বোচ্চ \(2^n\)টি
\textenglish{output}-এর মধ্যে একটি নির্দিষ্ট
\textenglish{output} সক্রিয় করে, তাকে
\textbf{ডিকোডার (\textenglish{Decoder})} বলে।
\par\endgroup


\[
\boxed{
n\text{টি input}
\longrightarrow
2^n\text{টি output}
}
\]

উদাহরণ:

\begin{itemize}[label={{\bulletfont\textbullet}}, leftmargin=*, nosep]
\item
\textenglish{2-to-4 Decoder}:
২টি \textenglish{input}, ৪টি \textenglish{output}।

\item
\textenglish{3-to-8 Decoder}:
৩টি \textenglish{input}, ৮টি \textenglish{output}।

\item
\textenglish{4-to-16 Decoder}:
৪টি \textenglish{input}, ১৬টি \textenglish{output}।

\end{itemize}

% =========================================================
% 3.12.3 Working Principle
% =========================================================

\subsection{ডিকোডারের গঠন ও কার্যপ্রণালি}

ডিকোডারের \textenglish{input}-এ একটি বাইনারি কোড দেওয়া হয়।
কোডটির মান অনুযায়ী \textenglish{output line}-গুলোর মধ্যে
একটি নির্দিষ্ট \textenglish{output} সক্রিয় হয়।

উদাহরণস্বরূপ, \textenglish{2-to-4 Decoder}-এ—

\[
AB=00\Rightarrow Y_0,
\qquad
AB=01\Rightarrow Y_1
\]

\[
AB=10\Rightarrow Y_2,
\qquad
AB=11\Rightarrow Y_3
\]

অর্থাৎ \textenglish{input code} পরিবর্তিত হলে সক্রিয়
\textenglish{output}-ও পরিবর্তিত হয়।

% =========================================================
% 3.12.4 Characteristics
% =========================================================

\subsection{ডিকোডারের বৈশিষ্ট্য
(\textenglish{Characteristics})}

\begin{itemize}[label={{\bulletfont\textbullet}}, leftmargin=*, nosep]

\item
ডিকোডার একটি সমবায় লজিক সার্কিট।

\item
\(n\)টি \textenglish{input}-এর জন্য সর্বোচ্চ
\(2^n\)টি \textenglish{output} থাকতে পারে।

\item
প্রতিটি বাইনারি \textenglish{input combination}-এর জন্য
একটি নির্দিষ্ট \textenglish{output} নির্বাচন করা যায়।

\item
সাধারণ \textenglish{Active-High Decoder}-এ নির্বাচিত
\textenglish{output} \(1\) হয় এবং অন্য
\textenglish{output}-গুলো \(0\) থাকে।

\item
অনেক ডিকোডারে \textenglish{Enable input} থাকে,
যা সার্কিটের কার্যক্রম চালু বা বন্ধ করতে ব্যবহৃত হয়।

\end{itemize}

% =========================================================
% 3.12.5 2-to-4 Decoder
% =========================================================

\subsection{\textenglish{2-to-4 Decoder}}

একটি \textenglish{2-to-4 Decoder}-এ দুটি
\textenglish{input} \(A,B\) এবং চারটি
\textenglish{output} \(Y_0,Y_1,Y_2,Y_3\) থাকে।

\[
2^2=4
\]

অতএব, ২-বিট \textenglish{input}-এর চারটি সম্ভাব্য সমন্বয়ের জন্য
চারটি পৃথক \textenglish{output} থাকে।

% =========================================================
% 2-to-4 Decoder: 3-column compact box
% Truth Table | Boolean Equations | Logic Circuit
% =========================================================

\par\begingroup


\begin{minipage}[t]{0.29\textwidth}
\centering

\textbf{সত্যক সারণি}\\[-1mm]
{\footnotesize \textenglish{Truth Table}}

\vspace{4mm}

{\renewcommand{\arraystretch}{1.24}

\begin{tabular}{|c|c|c|}
\hline
\rowcolor{ICTLightBlue}
\(A\) & \(B\) & \textbf{সক্রিয়} \\
\hline
0 & 0 & \(Y_0\) \\ \hline
0 & 1 & \(Y_1\) \\ \hline
1 & 0 & \(Y_2\) \\ \hline
1 & 1 & \(Y_3\) \\ \hline
\end{tabular}

}

\end{minipage}
\hfill
\begin{minipage}[t]{0.24\textwidth}
\centering

\textbf{বুলিয়ান সমীকরণ}\\[-1mm]
{\footnotesize \textenglish{Boolean Equations}}

\vspace{3mm}

\[
\begin{array}{c}
Y_0=\bar A\bar B\\[2.2mm]
Y_1=\bar A B\\[2.2mm]
Y_2=A\bar B\\[2.2mm]
Y_3=AB
\end{array}
\]

\end{minipage}
\hfill
\begin{minipage}[t]{0.43\textwidth}
\centering

\textbf{লজিক সার্কিট}\\[-1mm]
{\footnotesize \textenglish{Logic Circuit}}

\vspace{1mm}

\begin{chThreeDiagram}
\begin{tikzpicture}[
    line width=0.65pt,
    scale=0.66,
    transform shape
]

\coordinate (A) at (0,0);
\draw (0,4.1) -- (0,-3.8);
\node[above,font=\small] at (0,4.1) {\(A\)};

\node[
    not gate US,
    draw,
    scale=0.85
] (notA) at (0.85,3.6) {};

\draw (0,3.6) -- (notA.input);
\draw (notA.output) -- (1.55,3.6);
\draw (1.55,3.6) -- (1.55,-3.8);

\node[above,font=\small] at (1.55,3.6)
{\(\bar A\)};

\coordinate (Abar) at (1.55,0);

\coordinate (B) at (2.45,0);
\draw (2.45,4.1) -- (2.45,-3.8);
\node[above,font=\small] at (2.45,4.1) {\(B\)};

\node[
    not gate US,
    draw,
    scale=0.85
] (notB) at (3.30,3.6) {};

\draw (2.45,3.6) -- (notB.input);
\draw (notB.output) -- (4.00,3.6);
\draw (4.00,3.6) -- (4.00,-3.8);

\node[above,font=\small] at (4.00,3.6)
{\(\bar B\)};

\coordinate (Bbar) at (4.00,0);

\node[
    and gate US,
    draw,
    minimum width=1.55cm,
    minimum height=1.12cm
] (g0) at (7.3,2.7) {};

\node[
    and gate US,
    draw,
    minimum width=1.55cm,
    minimum height=1.12cm
] (g1) at (7.3,0.9) {};

\node[
    and gate US,
    draw,
    minimum width=1.55cm,
    minimum height=1.12cm
] (g2) at (7.3,-0.9) {};

\node[
    and gate US,
    draw,
    minimum width=1.55cm,
    minimum height=1.12cm
] (g3) at (7.3,-2.7) {};

\foreach \g/\sone/\stwo in {
    g0/Abar/Bbar,
    g1/Abar/B,
    g2/A/Bbar,
    g3/A/B
}{
    \coordinate (\g-up) at
    ($(\g.west)+(0,0.28)$);

    \coordinate (\g-dn) at
    ($(\g.west)+(0,-0.28)$);

    \draw
    (\sone |- \g-up)
    -- (\g-up);

    \draw
    (\stwo |- \g-dn)
    -- (\g-dn);

    \fill
    (\sone |- \g-up)
    circle (1.35pt);

    \fill
    (\stwo |- \g-dn)
    circle (1.35pt);
}

\draw
(g0.output)
-- ++(0.8,0)
node[right] {\(Y_0\)};

\draw
(g1.output)
-- ++(0.8,0)
node[right] {\(Y_1\)};

\draw
(g2.output)
-- ++(0.8,0)
node[right] {\(Y_2\)};

\draw
(g3.output)
-- ++(0.8,0)
node[right] {\(Y_3\)};

\end{tikzpicture}
\end{chThreeDiagram}

\end{minipage}

\par\endgroup

\par\begingroup

\textenglish{2-to-4 Decoder} বাস্তবায়নে সাধারণভাবে
২টি \textenglish{NOT Gate} এবং
৪টি ২-\textenglish{input AND Gate} ব্যবহৃত হয়।
\par\endgroup


% =========================================================
% 2-to-4 Example
% =========================================================

\subsubsection*{উদাহরণ}

একটি \textenglish{2-to-4 Decoder}-এ
\(A=1\) এবং \(B=0\) হলে—

\[
Y_0=\bar A\bar B
=0\cdot1
=0
\]

\[
Y_1=\bar A B
=0\cdot0
=0
\]

\[
Y_2=A\bar B
=1\cdot1
=1
\]

\[
Y_3=AB
=1\cdot0
=0
\]

অতএব,

\[
\boxed{Y_2=1}
\]

এবং অন্য সব \textenglish{output} \(0\)।

% =========================================================
% 3.12.6 3-to-8 Decoder
% =========================================================

\subsection{\textenglish{3-to-8 Decoder}}

একটি \textenglish{3-to-8 Decoder}-এ তিনটি
\textenglish{input} \(A,B,C\) এবং আটটি
\textenglish{output} \(Y_0\) থেকে \(Y_7\) থাকে।

\[
2^3=8
\]

অর্থাৎ ৩-বিট \textenglish{input}-এর আটটি সম্ভাব্য সমন্বয়ের জন্য
আটটি পৃথক \textenglish{output} থাকে।

% =========================================================
% 3-to-8: FIXED THREE-COLUMN BOX
% =========================================================

\par\begingroup


\begin{minipage}[t]{0.25\textwidth}
\centering

\textbf{সত্যক সারণি}\\[-1mm]
{\footnotesize \textenglish{Truth Table}}

\vspace{4mm}

{\renewcommand{\arraystretch}{1.24}

\begin{tabular}{|c|c|c|c|}
\hline
\rowcolor{ICTLightBlue}
\(A\) & \(B\) & \(C\) & \textbf{সক্রিয়} \\
\hline
0 & 0 & 0 & \(Y_0\) \\ \hline
0 & 0 & 1 & \(Y_1\) \\ \hline
0 & 1 & 0 & \(Y_2\) \\ \hline
0 & 1 & 1 & \(Y_3\) \\ \hline
1 & 0 & 0 & \(Y_4\) \\ \hline
1 & 0 & 1 & \(Y_5\) \\ \hline
1 & 1 & 0 & \(Y_6\) \\ \hline
1 & 1 & 1 & \(Y_7\) \\ \hline
\end{tabular}

}

\end{minipage}
\hfill
\begin{minipage}[t]{0.25\textwidth}
\centering

\textbf{বুলিয়ান সমীকরণ}\\[-1mm]
{\footnotesize \textenglish{Boolean Equations}}

\vspace{2mm}

\[
\begin{array}{c}
Y_0=\bar A\bar B\bar C\\[1.5mm]
Y_1=\bar A\bar B C\\[1.5mm]
Y_2=\bar A B\bar C\\[1.5mm]
Y_3=\bar A BC\\[1.5mm]
Y_4=A\bar B\bar C\\[1.5mm]
Y_5=A\bar B C\\[1.5mm]
Y_6=AB\bar C\\[1.5mm]
Y_7=ABC
\end{array}
\]

\end{minipage}
\hfill
\begin{minipage}[t]{0.46\textwidth}
\centering

\textbf{লজিক সার্কিট}\\[-1mm]
{\footnotesize \textenglish{Logic Circuit}}

\vspace{1mm}

\begin{chThreeDiagram}
\begin{tikzpicture}[
    line width=0.58pt,
    scale=0.51,
    transform shape
]

\coordinate (A) at (0,0);
\draw (0,7.9) -- (0,-7.5);
\node[above,font=\small] at (0,7.9) {\(A\)};

\node[
    not gate US,
    draw,
    scale=0.82
] (notA) at (0.8,7.35) {};

\draw (0,7.35) -- (notA.input);
\draw (notA.output) -- (1.45,7.35);
\draw (1.45,7.35) -- (1.45,-7.5);

\node[above,font=\small] at (1.45,7.35)
{\(\bar A\)};

\coordinate (Abar) at (1.45,0);

\coordinate (B) at (2.35,0);
\draw (2.35,7.9) -- (2.35,-7.5);
\node[above,font=\small] at (2.35,7.9) {\(B\)};

\node[
    not gate US,
    draw,
    scale=0.82
] (notB) at (3.15,7.35) {};

\draw (2.35,7.35) -- (notB.input);
\draw (notB.output) -- (3.80,7.35);
\draw (3.80,7.35) -- (3.80,-7.5);

\node[above,font=\small] at (3.80,7.35)
{\(\bar B\)};

\coordinate (Bbar) at (3.80,0);

\coordinate (C) at (4.70,0);
\draw (4.70,7.9) -- (4.70,-7.5);
\node[above,font=\small] at (4.70,7.9) {\(C\)};

\node[
    not gate US,
    draw,
    scale=0.82
] (notC) at (5.50,7.35) {};

\draw (4.70,7.35) -- (notC.input);
\draw (notC.output) -- (6.15,7.35);
\draw (6.15,7.35) -- (6.15,-7.5);

\node[above,font=\small] at (6.15,7.35)
{\(\bar C\)};

\coordinate (Cbar) at (6.15,0);

\node[
    and gate US,
    draw,
    minimum width=1.55cm,
    minimum height=1.12cm
] (g0) at (10.9,6.7) {};

\node[
    and gate US,
    draw,
    minimum width=1.55cm,
    minimum height=1.12cm
] (g1) at (10.9,4.8) {};

\node[
    and gate US,
    draw,
    minimum width=1.55cm,
    minimum height=1.12cm
] (g2) at (10.9,2.9) {};

\node[
    and gate US,
    draw,
    minimum width=1.55cm,
    minimum height=1.12cm
] (g3) at (10.9,1.0) {};

\node[
    and gate US,
    draw,
    minimum width=1.55cm,
    minimum height=1.12cm
] (g4) at (10.9,-1.0) {};

\node[
    and gate US,
    draw,
    minimum width=1.55cm,
    minimum height=1.12cm
] (g5) at (10.9,-2.9) {};

\node[
    and gate US,
    draw,
    minimum width=1.55cm,
    minimum height=1.12cm
] (g6) at (10.9,-4.8) {};

\node[
    and gate US,
    draw,
    minimum width=1.55cm,
    minimum height=1.12cm
] (g7) at (10.9,-6.7) {};

\newcommand{\decoderthreeinputs}[4]{%

    \coordinate (#1-u) at
    ($(#1.west)+(0,0.34)$);

    \coordinate (#1-m) at
    ($(#1.west)+(0,0)$);

    \coordinate (#1-d) at
    ($(#1.west)+(0,-0.34)$);

    \draw
    (#2 |- #1-u)
    -- (#1-u);

    \draw
    (#3 |- #1-m)
    -- (#1-m);

    \draw
    (#4 |- #1-d)
    -- (#1-d);

    \fill
    (#2 |- #1-u)
    circle (1.2pt);

    \fill
    (#3 |- #1-m)
    circle (1.2pt);

    \fill
    (#4 |- #1-d)
    circle (1.2pt);
}

\decoderthreeinputs{g0}{Abar}{Bbar}{Cbar}
\decoderthreeinputs{g1}{Abar}{Bbar}{C}
\decoderthreeinputs{g2}{Abar}{B}{Cbar}
\decoderthreeinputs{g3}{Abar}{B}{C}
\decoderthreeinputs{g4}{A}{Bbar}{Cbar}
\decoderthreeinputs{g5}{A}{Bbar}{C}
\decoderthreeinputs{g6}{A}{B}{Cbar}
\decoderthreeinputs{g7}{A}{B}{C}

\draw
(g0.output)
-- ++(0.75,0)
node[right] {\(Y_0\)};

\draw
(g1.output)
-- ++(0.75,0)
node[right] {\(Y_1\)};

\draw
(g2.output)
-- ++(0.75,0)
node[right] {\(Y_2\)};

\draw
(g3.output)
-- ++(0.75,0)
node[right] {\(Y_3\)};

\draw
(g4.output)
-- ++(0.75,0)
node[right] {\(Y_4\)};

\draw
(g5.output)
-- ++(0.75,0)
node[right] {\(Y_5\)};

\draw
(g6.output)
-- ++(0.75,0)
node[right] {\(Y_6\)};

\draw
(g7.output)
-- ++(0.75,0)
node[right] {\(Y_7\)};

\end{tikzpicture}
\end{chThreeDiagram}

\end{minipage}

\par\endgroup

\par\begingroup

\textenglish{3-to-8 Decoder} বাস্তবায়নে সাধারণভাবে
৩টি \textenglish{NOT Gate} এবং
৮টি ৩-\textenglish{input AND Gate} ব্যবহৃত হয়।
\par\endgroup


% =========================================================
% 3-to-8 Example
% =========================================================

\subsubsection*{উদাহরণ}

যদি

\[
ABC=101
\]

হয়, তবে

\[
(101)_2=(5)_{10}
\]

সুতরাং,

\[
\boxed{Y_5=1}
\]

এবং অন্য সব \textenglish{output} \(0\)।

সমীকরণ থেকেও—

\[
Y_5=A\bar B C
=1\cdot1\cdot1
=1
\]

% =========================================================
% 3.12.7 Applications
% =========================================================

\subsection{ডিকোডারের ব্যবহার
(\textenglish{Applications of Decoder})}
\begin{itemize}[label={{\bulletfont\textbullet}}, leftmargin=*, nosep]

\item
কম্পিউটার মেমোরির নির্দিষ্ট
\textenglish{address} বা অবস্থান নির্বাচন করতে।

\item
বাইনারি কোড অনুযায়ী নির্দিষ্ট
ডিভাইস বা সার্কিট সক্রিয় করতে।

\item
ডিজিটাল নিয়ন্ত্রণ ব্যবস্থায় নির্দিষ্ট
\textenglish{output} নির্বাচন করতে।

\item
প্রসেসর ও ডিজিটাল সিস্টেমে
\textenglish{instruction/address decoding}-এ।

\item
\textenglish{BCD-to-Seven-Segment Decoder}-এর মাধ্যমে
\textenglish{Seven-Segment Display} নিয়ন্ত্রণে।

\end{itemize}

% =========================================================
% 3.12.8 Relationship
% =========================================================

\subsection{এনকোডার ও ডিকোডারের সম্পর্ক
(\textenglish{Relationship between Encoder and Decoder})}

এনকোডার (\textenglish{Encoder}) ও ডিকোডার
(\textenglish{Decoder}) পরস্পর বিপরীতধর্মী কাজ করে।

এনকোডার সক্রিয় \textenglish{input signal}-এর অবস্থানকে
বাইনারি \textenglish{code}-এ রূপান্তর করে।
অন্যদিকে ডিকোডার সেই বাইনারি \textenglish{code}-এর মান অনুযায়ী
একটি নির্দিষ্ট \textenglish{output} নির্বাচন বা সক্রিয় করে।

কম্পিউটার ব্যবস্থায় মানুষের দেওয়া কিছু \textenglish{input}
প্রথমে উপযুক্ত কোডে রূপান্তরিত হয়ে প্রক্রিয়াকৃত হতে পারে।
পরবর্তীতে প্রয়োজন অনুযায়ী সেই কোডের ভিত্তিতে নির্দিষ্ট
\textenglish{output} নির্বাচন বা প্রদর্শনের কাজে
ডিকোডিং ব্যবহৃত হতে পারে।

\begin{center}

\begin{chThreeDiagram}
\begin{tikzpicture}[
    >=Stealth,
    line width=0.8pt,
    node distance=0.65cm
]

\node[
    draw,
    rounded corners=2pt,
    fill=ICTLightBlue,
    minimum width=2.1cm,
    minimum height=0.85cm,
    align=center
] (i) {\textenglish{Input Signal}};

\node[
    draw,
    rounded corners=2pt,
    fill=ICTLightGreen,
    minimum width=1.9cm,
    minimum height=0.85cm,
    right=of i
] (e) {\textenglish{Encoder}};

\node[
    draw,
    rounded corners=2pt,
    fill=ICTLightOrange,
    minimum width=2.1cm,
    minimum height=0.85cm,
    right=of e
] (c) {\textenglish{Binary Code}};

\node[
    draw,
    rounded corners=2pt,
    fill=ICTLightGreen,
    minimum width=1.9cm,
    minimum height=0.85cm,
    right=of c
] (d) {\textenglish{Decoder}};

\node[
    draw,
    rounded corners=2pt,
    fill=ICTLightBlue,
    minimum width=2.4cm,
    minimum height=0.85cm,
    right=of d
] (o) {\textenglish{Selected Output}};

\draw[->] (i) -- (e);
\draw[->] (e) -- (c);
\draw[->] (c) -- (d);
\draw[->] (d) -- (o);

\end{tikzpicture}
\end{chThreeDiagram}

\end{center}

\par\begingroup


\textbf{সহজভাবে মনে রাখবে:}\quad

\textenglish{Encoder: Signal}
\(\rightarrow\)
\textenglish{Code},

\qquad

\textenglish{Decoder: Code}
\(\rightarrow\)
\textenglish{Selected Output}।

ডিকোডারের মূল কাজ হলো
\textenglish{input code} অনুযায়ী
নির্দিষ্ট \textenglish{output} নির্বাচন করা।

\par\endgroup


% =========================================================
% 3.12.9 Comparison
% =========================================================

\subsection{এনকোডার ও ডিকোডারের তুলনামূলক পার্থক্য}

\begin{center}

\footnotesize

\renewcommand{\arraystretch}{1.08}

\begin{tabularx}
{\textwidth}
{|L{0.18\textwidth}|Y|Y|}

\hline

\rowcolor{ICTLightBlue}

\textbf{তুলনার বিষয়}
&
\textbf{এনকোডার (\textenglish{Encoder})}
&
\textbf{ডিকোডার (\textenglish{Decoder})}
\\
\hline

সংজ্ঞা
&
সক্রিয় \textenglish{input}-এর অবস্থানকে
নির্দিষ্ট কোডে রূপান্তর করে।
&
কোডেড \textenglish{input} অনুযায়ী
নির্দিষ্ট \textenglish{output} সক্রিয় করে।
\\
\hline

প্রধান কাজ
&
সংকেতকে কোডে রূপান্তর করা।
&
কোড অনুযায়ী \textenglish{output} নির্বাচন করা।
\\
\hline

\textenglish{Input} সংখ্যা
&
সাধারণত \(2^n\)টি।
&
সাধারণত \(n\)টি।
\\
\hline

\textenglish{Output} সংখ্যা
&
সাধারণত \(n\)টি।
&
সর্বোচ্চ \(2^n\)টি।
\\
\hline

উদাহরণ
&
\textenglish{4-to-2},
\textenglish{8-to-3}
&
\textenglish{2-to-4},
\textenglish{3-to-8}
\\
\hline

কার্যপ্রণালি
&
সক্রিয় \textenglish{input}-এর অবস্থানকে
বাইনারি কোডে প্রকাশ করে।
&
বাইনারি \textenglish{input}-এর মান অনুযায়ী
নির্দিষ্ট \textenglish{output} সক্রিয় করে।
\\
\hline

ব্যবহার
&
\textenglish{Input signal coding}
ও সংকেত শনাক্তকরণে।
&
মেমোরি নির্বাচন, ঠিকানা ডিকোডিং
ও \textenglish{output selection}-এ।
\\
\hline

\end{tabularx}

\end{center}

\par\begingroup


\textbf{মনে রাখার কৌশল:}
এনকোডার কোড তৈরি করে,
আর ডিকোডার সেই কোড অনুযায়ী
নির্দিষ্ট \textenglish{output} নির্বাচন করে।

\par\endgroup

\subsection{Enable, ঠিকানা নির্বাচন ও function তৈরি}
একটি আদর্শ active-high decoder-এ $E=1$ হলে selected output সক্রিয়; $E=0$ হলে সব আউটপুট শূন্য। যেমন 2-to-4 decoder-এ $Y_0=E\overline A\overline B$, $Y_1=E\overline AB$, $Y_2=EA\overline B$, $Y_3=EAB$। E নিয়ন্ত্রণ ইনপুট; এর জন্য address bits-এর সংখ্যা তিন হয়ে যায় না।

\chThreeExample{১}{একটি memory block নির্বাচন}
চারটি memory block-এর নির্বাচক $A,B$। $AB=10,E=1$।
\chThreeAnswer $Y_2=1$, অন্যগুলো ০। $E=0$ করলে $AB=10$ থাকলেও কোনো block সক্রিয় হবে না। বাস্তব active-low decoder-এ selected line শূন্য হতে পারে; এটি ভিন্ন polarity।

\chThreeExample{২}{ডিকোডার দিয়ে XOR}
\chThreeAnswer 2-to-4 decoder-এর $Y_1=\overline AB$ এবং $Y_2=A\overline B$ একটি OR-এ দিলে $F=Y_1+Y_2=A\oplus B$। তাই decoder-এর minterm outputs থেকে বুলিয়ান function গঠন করা যায়।

\chThreeExample{৩}{সাত segment-এ অঙ্ক দেখানো}
চার বিট BCD $0101$ ইনপুট দিলে display-এ ৫ দেখাতে সাতটি segment-এর নির্দিষ্ট কয়েকটি জ্বালাতে হয়।
\chThreeAnswer BCD-to-seven-segment decoder-এর সাতটি segment output $a,b,c,d,e,f,g$ থাকে; ৫-এর জন্য $a,c,d,f,g$ জ্বলে। এটি সাধারণ 4-to-16 line decoder নয়। display-এর common-anode/common-cathode ব্যবস্থা অনুযায়ী active polarity ভিন্ন হতে পারে।

\chThreeExample{৪}{ধারণাভিত্তিক MCQ}
সাধারণ binary decoder-এর পাঁচটি address input থাকলে কয়টি output line? (ক) ১০, (খ) ১৬, (গ) ২৫, (ঘ) ৩২।
\chThreeAnswer (ঘ) ৩২; $2^5$। সব line একসঙ্গে সক্রিয় হয় না।

% Mahbub teaching-type follow-up, 2026-10-14
\chThreeExample{৫}{সক্রিয় output দেখে decoder input উদ্ধার}
তিন-input, আট-output active-high decoder enabled আছে। শুধু $Y_6=1$। input কী এবং তার minterm কী?
\chThreeAnswer $6=(110)_2$, তাই $ABC=110$ এবং $Y_6=AB\overline C$। A=১, B=১, C=০ বসালে $Y_6=1$; অন্য output-গুলোর অন্তত একটি প্রয়োজনীয় factor শূন্য। active-low device-এ নির্বাচিত রেখার level ০ হতে পারে; রেখার পরিচয় ৬ থাকলেও level-এর নিয়ম বদলাবে।

